\documentclass[11pt,a4paper]{article}

% --- Essential Packages ---
\usepackage[utf8]{inputenc}
\usepackage[margin=1in]{geometry}
\usepackage{amsmath,amssymb,amsthm,mathtools}
\usepackage{xcolor}
\usepackage{tcolorbox}
\usepackage{fancyhdr}
\usepackage{enumitem}
\usepackage{titlesec}
\usepackage{hyperref}
\usepackage{bookmark}

% --- Color Palette ---
\definecolor{primary}{RGB}{20, 42, 68}       % Deep Navy
\definecolor{secondary}{RGB}{38, 80, 120}    % Slate Blue
\definecolor{accentteal}{RGB}{24, 112, 116}  % Dark Teal
\definecolor{accentgold}{RGB}{165, 105, 30}  % Warm Gold
\definecolor{darkgray}{RGB}{50, 50, 50}

% Box background tints
\definecolor{thmbg}{RGB}{244, 248, 253}
\definecolor{thmfr}{RGB}{38, 80, 120}

\definecolor{defbg}{RGB}{243, 249, 248}
\definecolor{deffr}{RGB}{24, 112, 116}

\definecolor{exbg}{RGB}{252, 250, 246}
\definecolor{exfr}{RGB}{165, 105, 30}

\definecolor{notebg}{RGB}{248, 249, 250}
\definecolor{notefr}{RGB}{110, 120, 130}

% --- Hyperref Setup ---
\hypersetup{
    colorlinks=true,
    linkcolor=secondary,
    citecolor=accentteal,
    urlcolor=secondary,
    pdfauthor={Aryan Maurya},
    pdftitle={Abstract Algebra: Group Actions, Permutations and Ring Theory},
    pdfsubject={Lecture Notes in Abstract Algebra},
    bookmarksnumbered=true
}

% --- Custom Environments using tcolorbox ---
\newtcolorbox{theorembox}[1][]{
    colback=thmbg,
    colframe=thmfr,
    coltitle=white,
    fonttitle=\bfseries,
    arc=2mm,
    boxrule=0.6pt,
    leftrule=3.5pt,
    top=6pt,
    bottom=6pt,
    left=8pt,
    right=8pt,
    title={#1}
}

\newtcolorbox{defbox}[1][]{
    colback=defbg,
    colframe=deffr,
    coltitle=white,
    fonttitle=\bfseries,
    arc=2mm,
    boxrule=0.6pt,
    leftrule=3.5pt,
    top=6pt,
    bottom=6pt,
    left=8pt,
    right=8pt,
    title={#1}
}

\newtcolorbox{exbox}[1][]{
    colback=exbg,
    colframe=exfr,
    coltitle=white,
    fonttitle=\bfseries,
    arc=2mm,
    boxrule=0.6pt,
    leftrule=3.5pt,
    top=6pt,
    bottom=6pt,
    left=8pt,
    right=8pt,
    title={#1}
}

\newtcolorbox{notebox}[1][]{
    colback=notebg,
    colframe=notefr,
    coltitle=white,
    fonttitle=\bfseries,
    arc=2mm,
    boxrule=0.6pt,
    leftrule=3pt,
    top=5pt,
    bottom=5pt,
    left=8pt,
    right=8pt,
    title={#1}
}

% --- Counter and Environment Definitions ---
\newcounter{thmcount}
\counterwithin{thmcount}{section}

\newenvironment{theorem}[1][]%
{\refstepcounter{thmcount}\begin{theorembox}[Theorem \thethmcount\ifx&#1&\else: #1\fi]}%
{\end{theorembox}}

\newenvironment{proposition}[1][]%
{\refstepcounter{thmcount}\begin{theorembox}[Proposition \thethmcount\ifx&#1&\else: #1\fi]}%
{\end{theorembox}}

\newenvironment{corollary}[1][]%
{\refstepcounter{thmcount}\begin{theorembox}[Corollary \thethmcount\ifx&#1&\else: #1\fi]}%
{\end{theorembox}}

\newenvironment{definition}[1][]%
{\refstepcounter{thmcount}\begin{defbox}[Definition \thethmcount\ifx&#1&\else: #1\fi]}%
{\end{defbox}}

\newenvironment{example}[1][]%
{\refstepcounter{thmcount}\begin{exbox}[Example \thethmcount\ifx&#1&\else: #1\fi]}%
{\end{exbox}}

\newenvironment{exercise}[1][]%
{\refstepcounter{thmcount}\begin{exbox}[Exercise \thethmcount\ifx&#1&\else: #1\fi]}%
{\end{exbox}}

\newenvironment{note}[1][]%
{\begin{notebox}[\textbf{Note}\ifx&#1&\else: #1\fi]}%
{\end{notebox}}

\newenvironment{remark}[1][]%
{\begin{notebox}[\textbf{Remark}\ifx&#1&\else: #1\fi]}%
{\end{notebox}}

% --- Proof Customization ---
\renewcommand{\qedsymbol}{$\blacksquare$}

% --- Header and Footer Styling ---
\pagestyle{fancy}
\fancyhf{}
\fancyhead[L]{\small\textcolor{secondary}{\nouppercase{\leftmark}}}
\fancyhead[R]{\small\textcolor{darkgray}{Abstract Algebra Notes}}
\fancyfoot[C]{\small\textcolor{darkgray}{\thepage}}
\renewcommand{\headrulewidth}{0.4pt}
\renewcommand{\footrulewidth}{0pt}

% Section Titles Styling
\titleformat{\section}
  {\normalfont\Large\bfseries\color{primary}}{\thesection}{1em}{}[\color{secondary}\titlerule]
\titleformat{\subsection}
  {\normalfont\large\bfseries\color{secondary}}{\thesubsection}{1em}{}

% --- Document Metadata ---
\title{
    \vspace{-1.5cm}
    {\normalsize\scshape\color{accentteal} Advanced Pure Mathematics \textbar\ Abstract Algebra}\\[0.4cm]
    {\Huge\bfseries\color{primary} Group Actions, Permutations \\ \& Introduction to Ring Theory}\\[0.3cm]
    {\large\color{darkgray} Master Calibrated Lecture Notes and Comprehensive Mathematical Exposition}
}
\author{\textbf{Aryan Maurya}}
\date{\today}

\begin{document}

\pagenumbering{roman}
\maketitle
\thispagestyle{empty}

\begin{center}
\begin{minipage}{0.92\textwidth}
\small\itshape
\centering
These notes provide a unified and rigorous exposition of group actions, conjugation, permutation representations, Cayley's theorem, orbit-stabilizer theory, the class equation, $p$-groups, conjugacy classes in symmetric groups, and introductory ring theory with matrix ring constructions.
\end{minipage}
\end{center}

\vspace{0.3cm}
\tableofcontents
\vspace{0.6cm}
\hrule
\vspace{0.8cm}

\newpage
\pagenumbering{arabic}

% =============================================================================
% PART I: GROUP ACTIONS & PERMUTATION REPRESENTATIONS
% =============================================================================
\part*{Part I: Group Actions and Permutation Representations}
\addcontentsline{toc}{part}{Part I: Group Actions and Permutation Representations}

\section{Group Actions: Left and Right Actions}

\begin{definition}[Left Group Action]
Let $G$ be a group and $A$ a non-empty set. A \textbf{left group action} of $G$ on $A$ is a mapping
\[
* : G \times A \longrightarrow A, \quad (g, a) \longmapsto g * a
\]
satisfying the following two axioms:
\begin{enumerate}[label=(\roman*)]
    \item $g_1 * (g_2 * a) = (g_1 g_2) * a \quad \forall g_1, g_2 \in G, \; \forall a \in A$,
    \item $e * a = a \quad \forall a \in A$, where $e$ is the identity element of $G$.
\end{enumerate}
This action is also referred to as an action on $A$ by left, or a \textbf{left action of $G$ on $A$}.
\end{definition}

\begin{definition}[Right Group Action]
We can also define an action of $G$ on $A$ from the right by considering a mapping
\[
* : A \times G \longrightarrow A, \quad (a, g) \longmapsto a * g
\]
satisfying:
\begin{enumerate}[label=(\roman*)]
    \item $(a * g_1) * g_2 = a * (g_1 g_2) \quad \forall g_1, g_2 \in G, \; \forall a \in A$,
    \item $a * e = a \quad \forall a \in A$.
\end{enumerate}
\end{definition}

\subsection{Fundamental Examples}

\begin{example}[Trivial $G$-Action]
Let $G$ be any group and $A \neq \emptyset$. Define
\[
* : G \times A \longrightarrow A \quad \text{by} \quad g * a = a \quad \forall g \in G, \; \forall a \in A.
\]
Let $g_1, g_2 \in G$ and $a \in A$. Then:
\begin{align*}
g_1 * (g_2 * a) &= g_1 * a = a, \\
(g_1 g_2) * a &= a.
\end{align*}
Thus $g_1 * (g_2 * a) = (g_1 g_2) * a$, and clearly $e * a = a$.
This shows that $*$ is a group action of $G$ on $A$, known as the \textbf{trivial $G$-action}.
\end{example}

\begin{example}[Left Regular Action / Translation by Left]
Take $A = G$. Define
\[
* : G \times A \longrightarrow A \quad \text{by} \quad g * a = ga \in A \quad \forall g \in G, \; \forall a \in A = G.
\]
Verification of axioms:
\begin{enumerate}[label=(\roman*)]
    \item For $g_1, g_2 \in G$ and $a \in A$:
    \[
    g_1 * (g_2 * a) = g_1 * (g_2 a) = g_1(g_2 a) = (g_1 g_2)a = (g_1 g_2) * a.
    \]
    \item $e * a = ea = a$.
\end{enumerate}
Therefore, $*$ is a group action of $G$ on $A (= G)$. This action of $G$ on itself is known as the \textbf{action by left} or \textbf{translation by left}.
\end{example}

\begin{example}[Action by Right Translation]
Consider again $A = G$ and a mapping $* : G \times A \to A$.
\begin{itemize}
    \item If we define $g * a = ag$, then for $g_1, g_2 \in G$ and $a \in A$:
    \[
    g_1 * (g_2 * a) = g_1 * (ag_2) = (ag_2)g_1 = a(g_2 g_1) \neq a(g_1 g_2) = (g_1 g_2) * a.
    \]
    So this does \emph{not} define a left group action in general.
    \item However, if we define $* : G \times A \to A$ as
    \[
    g * a = ag^{-1} \quad \forall g \in G, \; \forall a \in A = G,
    \]
    then for all $g_1, g_2 \in G$ and $a \in A$:
    \begin{align*}
    g_1 * (g_2 * a) &= g_1 * (ag_2^{-1}) \\
    &= (ag_2^{-1})g_1^{-1} \\
    &= a(g_2^{-1} g_1^{-1}) \\
    &= a(g_1 g_2)^{-1} \\
    &= (g_1 g_2) * a,
    \end{align*}
    and
    \[
    e * a = ae^{-1} = a.
    \]
    Therefore, $*$ is a group action, referred to as the \textbf{action of $G$ by right translation} (or right multiplication). This is also known as the \textbf{regular action of $G$ on itself}.
\end{itemize}
\end{example}

\section{Actions by Conjugation}

\begin{example}[Action by Conjugation on $G$]
Let $G$ be a group. Take $A = G$. Define
\[
* : G \times A \longrightarrow A \quad \text{by} \quad g * a = gag^{-1} \quad \forall g \in G, \; \forall a \in A.
\]
Then $*$ is a group action on $G$, known as the \textbf{action by conjugation}.
\begin{proof}[Verification of axioms]
Take $g_1, g_2 \in G$ and $a \in A (= G)$:
\begin{enumerate}[label=(\roman*)]
    \item 
    \begin{align*}
    g_1 * (g_2 * a) &= g_1 * (g_2 a g_2^{-1}) \\
    &= g_1 g_2 a g_2^{-1} g_1^{-1} \\
    &= g_1 g_2 a (g_1 g_2)^{-1} \\
    &= (g_1 g_2) * a.
    \end{align*}
    \item 
    \[
    e * a = eae^{-1} = a.
    \]
\end{enumerate}
Therefore, the axioms hold and $*$ is a valid group action.
\end{proof}
\end{example}

\begin{example}[Action by Conjugation on Cosets]
Let $G$ be a group and $H \trianglelefteq G$ be a normal subgroup. Take
\[
A = G/H = \{aH \mid a \in G\},
\]
the set of all left cosets of $H$ in $G$. Define
\[
* : G \times A \longrightarrow A \quad \text{by} \quad g * (xH) = gxg^{-1}H \quad \forall g \in G, \; \forall xH \in A.
\]
\begin{proof}[Verification of axioms]
Take $g_1, g_2 \in G$ and $xH \in A$:
\begin{enumerate}[label=(\roman*)]
    \item 
    \begin{align*}
    g_1 * (g_2 * xH) &= g_1 * (g_2 x g_2^{-1}H) \\
    &= g_1 g_2 x g_2^{-1} g_1^{-1}H \\
    &= g_1 g_2 x (g_1 g_2)^{-1}H \\
    &= (g_1 g_2) * xH.
    \end{align*}
    \item 
    \[
    e * xH = exe^{-1}H = xH.
    \]
\end{enumerate}
Therefore, $*$ is a group action.
\end{proof}
\end{example}

\begin{example}[Action by Conjugation on Subgroups]
Let $G$ be a group and let $A$ be the set of all subgroups of $G$. Define
\[
* : G \times A \longrightarrow A \quad \text{by} \quad g * H = gHg^{-1} \quad \forall g \in G, \; \forall H \in A.
\]
Then $*$ is also a group action.
\begin{proof}[Verification of axioms]
For $g_1, g_2 \in G$ and $H \in A$:
\begin{enumerate}[label=(\roman*)]
    \item 
    \begin{align*}
    g_1 * (g_2 * H) &= g_1 * (g_2 H g_2^{-1}) \\
    &= g_1 g_2 H g_2^{-1} g_1^{-1} \\
    &= g_1 g_2 H (g_1 g_2)^{-1} \\
    &= (g_1 g_2) * H.
    \end{align*}
    \item 
    \[
    e * H = e H e^{-1} = H.
    \]
\end{enumerate}
Therefore, $*$ is a group action on the set of subgroups.
\end{proof}
\end{example}

\begin{remark}
If $* : G \times A \to A$ is an action, then because $*$ is a map, if $x, y \in A$ such that $x = y$, then for all $g \in G$:
\[
(g, x) = (g, y) \implies g * x = g * y.
\]
\end{remark}

\section{Group Actions and Permutation Representations}

\begin{theorem}
Let $G$ be any group and $A$ be a non-empty set. Then any homomorphism from $G$ to $\operatorname{Sym}(A)$, the symmetric group on $A$, defines an action of $G$ on $A$.

Conversely, every action of $G$ on $A$ induces a homomorphism from $G$ to $\operatorname{Sym}(A)$.
\end{theorem}

\begin{proof}
We establish the theorem in two parts:

\paragraph{Part 1: Homomorphism induces an action.}
Let $F : G \to \operatorname{Sym}(A)$ be a group homomorphism. For each $g \in G$, denote $F(g) = \sigma_g$, where $\sigma_g : A \to A$ is a bijective map. Since $F$ is a homomorphism, for any $g_1, g_2 \in G$:
\[
F(g_1 g_2) = F(g_1) \circ F(g_2) \implies \sigma_{g_1 g_2} = \sigma_{g_1} \circ \sigma_{g_2}.
\]
Define $* : G \times A \to A$ by:
\[
g * a = \sigma_g(a) \quad \forall g \in G, \; \forall a \in A.
\]
Then $*$ is a map. For any $g_1, g_2 \in G$ and $a \in A$:
\begin{align*}
g_1 * (g_2 * a) &= g_1 * (\sigma_{g_2}(a)) \\
&= \sigma_{g_1}(\sigma_{g_2}(a)) \\
&= (\sigma_{g_1} \circ \sigma_{g_2})(a) \\
&= \sigma_{g_1 g_2}(a) \\
&= (g_1 g_2) * a.
\end{align*}
Furthermore, since $F$ is a homomorphism, $F(e) = \sigma_e = I_A$, where $I_A$ is the identity permutation on $A$. Thus:
\[
e * a = \sigma_e(a) = I_A(a) = a.
\]
Therefore, $*$ is a group action of $G$ on $A$.

\paragraph{Part 2: Action induces a homomorphism.}
Conversely, suppose $* : G \times A \to A$ is a group action. Define $F : G \to \operatorname{Sym}(A)$ by setting $F(g) = \sigma_g$ for each $g \in G$, where $\sigma_g : A \to A$ is defined as:
\[
\sigma_g(a) = g * a \quad \forall a \in A.
\]
First, we show that $\sigma_g$ is a bijective map from $A$ to $A$:
\begin{itemize}
    \item \textbf{One-to-one (Injective):} Let $a_1, a_2 \in A$ such that $\sigma_g(a_1) = \sigma_g(a_2)$. Then:
    \begin{align*}
    g * a_1 = g * a_2 &\implies g^{-1} * (g * a_1) = g^{-1} * (g * a_2) \\
    &\implies (g^{-1} g) * a_1 = (g^{-1} g) * a_2 \quad (\text{as } * \text{ is an action}) \\
    &\implies e * a_1 = e * a_2 \\
    &\implies a_1 = a_2.
    \end{align*}
    Hence, $\sigma_g$ is one-to-one.
    
    \item \textbf{Onto (Surjective):} Let $a \in A$. Then $g^{-1} * a \in A$, and:
    \[
    \sigma_g(g^{-1} * a) = g * (g^{-1} * a) = (g g^{-1}) * a = e * a = a.
    \]
    Hence, $\sigma_g$ is onto.
\end{itemize}
Thus $\sigma_g$ is bijective, meaning $\sigma_g \in \operatorname{Sym}(A)$, and $F$ is a well-defined map.

Next, we show that $F$ is a group homomorphism. For any $g_1, g_2 \in G$ and $a \in A$:
\begin{align*}
\sigma_{g_1 g_2}(a) &= (g_1 g_2) * a \\
&= g_1 * (g_2 * a) \\
&= g_1 * \sigma_{g_2}(a) \\
&= \sigma_{g_1}(\sigma_{g_2}(a)) \\
&= (\sigma_{g_1} \circ \sigma_{g_2})(a).
\end{align*}
Therefore, $\sigma_{g_1 g_2} = \sigma_{g_1} \circ \sigma_{g_2}$, which implies:
\[
F(g_1 g_2) = F(g_1) \circ F(g_2).
\]
Hence, $F$ is a group homomorphism.
\end{proof}

\begin{remark}
For any group action of a group $G$ on a set $A$, the induced homomorphism
\[
F : G \longrightarrow \operatorname{Sym}(A)
\]
is called the \textbf{associated permutation representation} of the given action.
\end{remark}

\begin{corollary}[Cayley's Theorem]
Every group is isomorphic to a subgroup of its permutation group (symmetric group).
\end{corollary}

\begin{proof}
Let $G$ be any group. We know that $G$ acts on itself ($A = G$) under the left regular action:
\[
* : G \times A \longrightarrow A, \quad g * a = ga \quad \forall g \in G, \; \forall a \in A = G.
\]
By the previous theorem, the corresponding permutation representation is the homomorphism:
\[
F : G \longrightarrow \operatorname{Sym}(G), \quad F(g) = \sigma_g, \quad \text{where } \sigma_g(a) = ga.
\]
We show that $F$ is one-to-one:
\begin{align*}
F(g_1) = F(g_2) &\implies \sigma_{g_1} = \sigma_{g_2} \\
&\implies \sigma_{g_1}(a) = \sigma_{g_2}(a) \quad \forall a \in G \\
&\implies g_1 a = g_2 a \quad \forall a \in G.
\end{align*}
Choosing $a = e$:
\[
g_1 e = g_2 e \implies g_1 = g_2.
\]
Therefore, $F$ is one-to-one, which means:
\[
\ker F = \{e\}.
\]
By the Fundamental Theorem of Group Homomorphisms (First Isomorphism Theorem):
\[
\frac{G}{\ker F} \cong F(G) \implies \frac{G}{\{e\}} \cong F(G) \implies G \cong F(G).
\]
Since $F(G) \le \operatorname{Sym}(G)$, every group $G$ is isomorphic to a subgroup of $\operatorname{Sym}(G)$.
\end{proof}

\section{Kernel of an Action}

\begin{definition}[Kernel of an Action]
Let $* : G \times A \to A$ be a group action. The \textbf{kernel} of $*$ is defined as:
\[
\ker * = \{g \in G \mid g * a = a, \; \forall a \in A\}.
\]
In other words, $\ker *$ consists of all elements of $G$ that fix every element of $A$.
\end{definition}

\begin{theorem}
Let $* : G \times A \to A$ be a group action. Then:
\[
\ker * \le G.
\]
\end{theorem}

\begin{proof}
We check the subgroup criteria:
\begin{enumerate}[label=(\roman*)]
    \item \textbf{Non-emptiness:} Since $e * a = a$ for all $a \in A$, we have $e \in \ker *$, so $\ker * \neq \emptyset$.
    
    \item \textbf{Closure under multiplication:} Let $g_1, g_2 \in \ker *$. Then:
    \[
    g_1 * a = a \quad \text{and} \quad g_2 * a = a \quad \forall a \in A.
    \]
    Then for any $a \in A$:
    \[
    (g_1 g_2) * a = g_1 * (g_2 * a) = g_1 * a = a.
    \]
    Thus $g_1 g_2 \in \ker *$.
    
    \item \textbf{Closure under inversion:} Let $g_1 \in \ker *$, so $g_1 * a = a$ for all $a \in A$. Then:
    \begin{align*}
    g_1^{-1} * (g_1 * a) = g_1^{-1} * a &\implies (g_1^{-1} g_1) * a = g_1^{-1} * a \\
    &\implies e * a = g_1^{-1} * a \\
    &\implies a = g_1^{-1} * a \quad \forall a \in A.
    \end{align*}
    Thus $g_1^{-1} \in \ker *$.
\end{enumerate}
Therefore, $\ker * \le G$.
\end{proof}

\begin{theorem}
If $F : G \to \operatorname{Sym}(A)$ is the permutation representation corresponding to the group action $*$, then:
\[
\ker F = \ker *.
\]
\end{theorem}

\begin{proof}
Recall that $F(g) = \sigma_g$ where $\sigma_g(a) = g * a$ for all $a \in A$. By the definition of the kernel of a group homomorphism:
\begin{align*}
\ker F &= \{g \in G \mid F(g) = I_A\} \\
&= \{g \in G \mid \sigma_g = \sigma_e\} \\
&= \{g \in G \mid \sigma_g(a) = \sigma_e(a) \quad \forall a \in A\} \\
&= \{g \in G \mid g * a = e * a \quad \forall a \in A\} \\
&= \{g \in G \mid g * a = a \quad \forall a \in A\} \\
&= \ker *.
\end{align*}
Hence, $\ker F = \ker *$.
\end{proof}

\section{Orbits and Stabilizers}

\begin{definition}[Stabilizer]
Let $*$ be a group action of a group $G$ on a set $A$, and let $a \in A$. The set
\[
G_a = \{g \in G \mid g * a = a\}
\]
is called the \textbf{stabilizer} (or isotropy group) of $a$ in $G$.
\end{definition}

\begin{theorem}
For any $a \in A$, the stabilizer $G_a$ is a subgroup of $G$:
\[
G_a \le G.
\]
\end{theorem}

\begin{proof}
We check the subgroup criteria:
\begin{enumerate}[label=(\roman*)]
    \item Since $e * a = a$, we have $e \in G_a$, so $G_a \neq \emptyset$.
    
    \item Let $g_1, g_2 \in G_a$. Then:
    \[
    g_1 * a = a \quad \text{and} \quad g_2 * a = a.
    \]
    Now, using the fact that $*$ is a group action:
    \[
    (g_1 g_2) * a = g_1 * (g_2 * a) = g_1 * a = a.
    \]
    Therefore:
    \[
    g_1 g_2 \in G_a \longrightarrow (1)
    \]
    
    \item Also, since $g_1 * a = a$:
    \begin{align*}
    g_1^{-1} * (g_1 * a) = g_1^{-1} * a &\implies (g_1^{-1} g_1) * a = g_1^{-1} * a \\
    &\implies e * a = g_1^{-1} * a \\
    &\implies g_1^{-1} * a = a.
    \end{align*}
    Therefore:
    \[
    g_1^{-1} \in G_a \longrightarrow (2)
    \]
\end{enumerate}
From (1) and (2), we conclude that $G_a \le G$.
\end{proof}

\begin{definition}[Orbit]
Let $*$ be a group action of $G$ on $A$, and let $a \in A$. The \textbf{orbit} of $a$ under $G$ is the set:
\[
G^a = \{g * a \mid g \in G\} \subseteq A.
\]
\end{definition}

\begin{theorem}[Orbit-Stabilizer Theorem]
Let $*$ be an action of a group $G$ on a set $A$, and let $a \in A$. There is a bijection between the orbit $G^a$ and the set of left cosets of the stabilizer, $G/G_a$.
\end{theorem}

\begin{proof}
Define the map
\[
\phi : G^a \longrightarrow G/G_a \quad \text{by} \quad \phi(g * a) = g G_a.
\]
We show that $\phi$ is well-defined, injective, and surjective:

\begin{itemize}
    \item \textbf{Well-defined:} Let $g_1 * a = g_2 * a$. Then:
    \begin{align*}
    g_2^{-1} * (g_1 * a) = g_2^{-1} * (g_2 * a) &\implies (g_2^{-1} g_1) * a = (g_2^{-1} g_2) * a \\
    &\implies (g_2^{-1} g_1) * a = e * a = a \\
    &\implies g_2^{-1} g_1 \in G_a \\
    &\implies g_1 G_a = g_2 G_a \\
    &\implies \phi(g_1 * a) = \phi(g_2 * a).
    \end{align*}
    Therefore, $\phi$ is well-defined.
    
    \item \textbf{One-to-one (Injective):} Let $\phi(g_1 * a) = \phi(g_2 * a)$. Then:
    \begin{align*}
    g_1 G_a = g_2 G_a &\implies g_2^{-1} g_1 \in G_a \\
    &\implies (g_2^{-1} g_1) * a = a \\
    &\implies g_2 * ((g_2^{-1} g_1) * a) = g_2 * a \\
    &\implies (g_2 g_2^{-1} g_1) * a = g_2 * a \\
    &\implies g_1 * a = g_2 * a.
    \end{align*}
    Therefore, $\phi$ is one-to-one.
    
    \item \textbf{Onto (Surjective):} Let $t \in G/G_a$. Then $t = g G_a$ for some $g \in G$. For this element $g$, $g * a \in G^a$ satisfies:
    \[
    \phi(g * a) = g G_a = t.
    \]
    Therefore, $\phi$ is onto.
\end{itemize}
Thus, $\phi$ is a bijection between $G^a$ and $G/G_a$.
\end{proof}

\begin{note}
If $G$ is a finite group, then:
\[
|G^a| = |G / G_a| = \frac{|G|}{|G_a|},
\]
which yields the fundamental relation:
\[
|G| = |G^a| \cdot |G_a|.
\]
\end{note}

\section{Orbits as Equivalence Classes and Conjugacy Classes}

\subsection{Equivalence Relation Induced by an Action}

\begin{theorem}
Let $*$ be a group action of a group $G$ on a set $A$, $* : G \times A \to A$. Define a relation $\sim$ on $A$ by:
\[
a \sim b \iff a = g * b \quad \text{for some } g \in G.
\]
Then $\sim$ is an equivalence relation on $A$.
\end{theorem}

\begin{proof}
We verify the three equivalence relation properties:
\begin{enumerate}[label=(\roman*)]
    \item \textbf{Reflexive:} For all $a \in A$, since $e \in G$ is the identity:
    \[
    a = e * a \implies a \sim a.
    \]
    Therefore, $\sim$ is reflexive.
    
    \item \textbf{Symmetric:} Let $a, b \in A$ such that $a \sim b$. Then there exists $g \in G$ such that:
    \[
    a = g * b.
    \]
    Applying $g^{-1} \in G$ to both sides:
    \begin{align*}
    g^{-1} * a &= g^{-1} * (g * b) \\
    &= (g^{-1} g) * b \\
    &= e * b = b.
    \end{align*}
    Since $g^{-1} \in G$, we have $b \sim a$. Therefore, $\sim$ is symmetric.
    
    \item \textbf{Transitive:} Let $a, b, c \in A$ such that $a \sim b$ and $b \sim c$. Then there exist $g_1, g_2 \in G$ such that:
    \[
    a = g_1 * b \quad \text{and} \quad b = g_2 * c.
    \]
    Substituting $b$ into the expression for $a$:
    \[
    a = g_1 * (g_2 * c) = (g_1 g_2) * c.
    \]
    Since $g_1 g_2 \in G$, we have $a \sim c$. Therefore, $\sim$ is transitive.
\end{enumerate}
Hence, $\sim$ is an equivalence relation on $A$.
\end{proof}

\subsection{Partition of the Set into Orbits}

Since $\sim$ is an equivalence relation on $A$, it partitions the set $A$ into its equivalence classes. For any $a \in A$, the equivalence class of $a$ is:
\begin{align*}
\operatorname{cl}(a) &= \{b \in A \mid a \sim b\} \\
&= \{b \in A \mid b = g * a \text{ for some } g \in G\} \\
&= \{g * a \mid g \in G\} \\
&= G^a \quad \text{(the orbit of $a$ under $G$)}.
\end{align*}
Consequently, the set $A$ decomposes as a disjoint union of its orbits:
\[
A = \bigcup_{a \in A} \operatorname{cl}(a) = \bigcup_{a \in A} G^a.
\]

\subsection{Conjugacy Classes}

If $G$ acts on itself ($A = G$) by conjugation, where $g * a = gag^{-1}$:
\begin{itemize}
    \item The stabilizer of $a$ is the \textbf{normalizer} (or centralizer) of $a$:
    \[
    G_a = N(a) = \{g \in G \mid gag^{-1} = a\} = \{g \in G \mid ga = ag\}.
    \]
    \item The orbit of $a$ is the \textbf{conjugacy class} of $a$:
    \[
    G^a = \operatorname{cl}(a) = \{g * a \mid g \in G\} = \{gag^{-1} \mid g \in G\}.
    \]
\end{itemize}

\begin{exercise}
Let $G$ be a group and $a, b \in G$. Then $a$ and $b$ are conjugate if there exists $c \in G$ such that $b = cac^{-1}$. Prove that this conjugacy relation $\sim$ is an equivalence relation on $G$.
\end{exercise}

\begin{proof}[Solution]
We check the three properties:
\begin{enumerate}[label=(\roman*)]
    \item \textbf{Reflexive:} Let $a \in G$. Since $e \in G$, we have:
    \[
    a = e a e^{-1}.
    \]
    Therefore, $a \sim a$, so $\sim$ is reflexive.
    
    \item \textbf{Symmetric:} Let $a, b \in G$ such that $a \sim b$. Then there exists $c \in G$ such that:
    \[
    a = c b c^{-1}.
    \]
    Multiplying on the left by $c^{-1}$ and on the right by $c$:
    \begin{align*}
    c^{-1} a &= c^{-1}(c b c^{-1}) = b c^{-1}, \\
    c^{-1} a c &= b c^{-1} c = b e = b.
    \end{align*}
    Writing $c = (c^{-1})^{-1}$, this gives:
    \[
    c^{-1} a (c^{-1})^{-1} = b.
    \]
    Since $c^{-1} \in G$, this shows $b \sim a$. Therefore, $\sim$ is symmetric.
    
    \item \textbf{Transitive:} Let $a, b, c \in G$ such that $a \sim b$ and $b \sim c$. Then there exist $d_1, d_2 \in G$ such that:
    \[
    a = d_1 b d_1^{-1} \quad \text{and} \quad b = d_2 c d_2^{-1}.
    \]
    Substituting $b$ into the first equation:
    \begin{align*}
    a &= d_1 (d_2 c d_2^{-1}) d_1^{-1} \\
    &= d_1 d_2 c d_2^{-1} d_1^{-1} \quad \{d_2^{-1} d_1^{-1} = (d_1 d_2)^{-1}\} \\
    &= (d_1 d_2) c (d_1 d_2)^{-1}.
    \end{align*}
    Since $d_1, d_2 \in G \implies d_1 d_2 \in G$, it follows that $a \sim c$.
\end{enumerate}
Therefore, $\sim$ is transitive, and hence an equivalence relation on $G$.
\end{proof}

\section{The Class Equation and Center of \texorpdfstring{$p$-Groups}{p-Groups}}

\subsection{Orbit-Stabilizer Applied to Conjugation}

By the Orbit-Stabilizer Theorem, for any $a \in G$, there exists a bijection between the conjugacy class $\operatorname{cl}(a)$ and the set of cosets $G/N(a)$, for all $a \in G$. Consequently:
\[
G = \bigcup_{a \in G} \operatorname{cl}(a).
\]
If $G$ is a finite group, then:
\begin{align*}
|\operatorname{cl}(a)| &= \left|\frac{G}{N(a)}\right| = \frac{|G|}{|N(a)|}, \\
|G| &= \sum_{a \in G} |\operatorname{cl}(a)|.
\end{align*}

\begin{proposition}
For any element $a \in G$:
\[
\operatorname{cl}(a) = \{a\} \iff a \in Z(G),
\]
where $Z(G) = \{z \in G \mid zg = gz \; \forall g \in G\}$ is the center of $G$.
\end{proposition}

\begin{proof}
By definition of the conjugacy class:
\begin{align*}
\operatorname{cl}(a) = \{a\} &\iff gag^{-1} = a \quad \forall g \in G \\
&\iff ga = ag \quad \forall g \in G \\
&\iff a \in Z(G).
\end{align*}
\end{proof}

\subsection{The Class Equation}

Separating the central elements (for which $|\operatorname{cl}(a)| = 1$) from the non-central elements yields the \textbf{Class Equation}:
\[
|G| = |Z(G)| + \sum_{a \notin Z(G)} \left|\frac{G}{N(a)}\right| \longrightarrow (1)
\]

\begin{theorem}
If $|G| = p^n$, where $p$ is a prime and $n \ge 1$, then:
\[
|Z(G)| \ge p.
\]
In other words, every non-trivial finite $p$-group has a non-trivial center.
\end{theorem}

\begin{proof}
Let $|Z(G)| = r$. By the class equation (1):
\[
|G| = |Z(G)| + \sum_{a \notin Z(G)} \left|\frac{G}{N(a)}\right|.
\]
For any $a \notin Z(G)$, $N(a) \neq G$, meaning $N(a)$ is a proper subgroup of $G$. By Lagrange's Theorem, $|N(a)|$ divides $|G| = p^n$, so:
\[
|N(a)| = p^{\alpha_a} \quad \text{with } \alpha_a < n.
\]
Then:
\[
\left|\frac{G}{N(a)}\right| = \frac{|G|}{|N(a)|} = \frac{p^n}{p^{\alpha_a}} = p^{n - \alpha_a}, \quad \text{where } n - \alpha_a > 0.
\]
This implies:
\[
p \mid p^{n - \alpha_a} \implies p \mid \left|\frac{G}{N(a)}\right| \implies p \mid \sum_{a \notin Z(G)} \left|\frac{G}{N(a)}\right|.
\]
Since $|G| = p^n$, we also have $p \mid |G|$. Rewriting the class equation:
\[
|Z(G)| = |G| - \sum_{a \notin Z(G)} \left|\frac{G}{N(a)}\right|.
\]
Since $p$ divides both terms on the right-hand side:
\[
p \mid |Z(G)| \implies p \mid r.
\]
Since the identity element $e \in Z(G)$, the center is never empty, so $|Z(G)| \ge 1$. Because $p$ divides $r$, we must have:
\[
r \ge p \implies |Z(G)| \ge p.
\]
\end{proof}

\begin{theorem}
A group of order $p^2$, where $p$ is prime, is abelian.
\end{theorem}

\begin{proof}
Let $|G| = p^2$. By the previous theorem:
\[
|Z(G)| \ge p.
\]
By Lagrange's Theorem, $|Z(G)|$ must divide $|G| = p^2$. Therefore, the only possibilities are:
\[
|Z(G)| = p \quad \text{or} \quad |Z(G)| = p^2.
\]
If $|Z(G)| = p$, then:
\[
\frac{|G|}{|Z(G)|} = \frac{p^2}{p} = p.
\]
Since any group of prime order is cyclic, the quotient group $G/Z(G)$ is cyclic. But it is a standard result that:
\[
G/Z(G) \text{ is cyclic} \implies G \text{ is abelian}.
\]
If $G$ is abelian, every element commutes with all other elements, so:
\[
G = Z(G) \implies |G| = |Z(G)| = p^2,
\]
which contradicts the assumption that $|Z(G)| = p$.

Therefore:
\[
|Z(G)| \neq p \quad \text{(contradiction)}.
\]
So we must have:
\[
|Z(G)| = p^2 = |G| \implies G = Z(G).
\]
Hence, $G$ is abelian.
\end{proof}

\begin{remark}
If $H \le G$ such that $H \subseteq Z(G)$, then $H \trianglelefteq G$ (any subgroup contained in the center is normal in $G$).
\end{remark}

\section{Conjugacy Classes in the Symmetric Group \texorpdfstring{$S_n$}{S\_n}}

\begin{theorem}[Number of Conjugates in $S_n$]
Let $\sigma$ be a permutation in $S_n$ and let $m_1, m_2, \dots, m_s$ be the distinct integers which appear in the cycle type of $\sigma$ (including $1$-cycles). For each $i \in \{1, 2, \dots, s\}$, assume $\sigma$ has $k_i$ cycles of length $m_i$ so that:
\[
\sum_{i=1}^s k_i m_i = n.
\]
Then the number of conjugates of $\sigma$ in $S_n$ (i.e., the size of the conjugacy class $|\operatorname{cl}(\sigma)|$) is:
\[
\frac{n!}{(k_1! \, m_1^{k_1})(k_2! \, m_2^{k_2}) \cdots (k_s! \, m_s^{k_s})}.
\]
\end{theorem}

\subsection{Examples in \texorpdfstring{$S_4$}{S\_4} and \texorpdfstring{$S_5$}{S\_5}}

\begin{example}[Conjugacy Classes in $S_4$]
\begin{enumerate}[label=(\alph*)]
    \item \textbf{Number of cycles of length $2$ in $S_4$:}
    A transposition in $S_4$ has cycle decomposition of the form $(a\,b)(c)(d)$. Here:
    \begin{itemize}
        \item $k_1 = 1$ cycle of length $m_1 = 2$,
        \item $k_2 = 2$ cycles of length $m_2 = 1$.
    \end{itemize}
    Applying the formula:
    \[
    \frac{4!}{(1! \cdot 2^1)(2! \cdot 1^2)} = \frac{4 \times 3 \times 2 \times 1}{2 \times 2} = 6.
    \]
    Explicitly, the conjugacy class of $(1\,2)$ consists of all $6$ transpositions:
    \[
    \operatorname{cl}((1\,2)) = \{(1\,2), (1\,3), (1\,4), (2\,3), (2\,4), (3\,4)\}.
    \]
    
    \item \textbf{Cycles of length $3$ in $S_4$:}
    A $3$-cycle has cycle type $(a\,b\,c)(d)$, so $k_1 = 1, m_1 = 3$ and $k_2 = 1, m_2 = 1$:
    \[
    |\operatorname{cl}((1\,2\,3))| = \frac{4!}{(1! \cdot 3^1)(1! \cdot 1^1)} = \frac{24}{3} = 8.
    \]
\end{enumerate}
\end{example}

\begin{example}[Conjugacy Classes in $S_5$]
\begin{enumerate}[label=(\alph*)]
    \item \textbf{Cycles of length $2$ in $S_5$:}
    A single transposition has cycle structure $(a\,b)(c)(d)(e)$, so $k_1 = 1, m_1 = 2$ and $k_2 = 3, m_2 = 1$:
    \[
    \text{No. of cycles of length } 2 = \frac{5!}{(3! \cdot 1^3)(1! \cdot 2^1)} = \frac{5 \times 4 \times 3 \times 2 \times 1}{(3 \times 2 \times 1) \times 2} = \frac{20}{2} = 10.
    \]
    
    \item \textbf{Permutations with cycle type $(2, 3)$ in $S_5$:}
    A product of a disjoint transposition and $3$-cycle has form $(a\,b)(c\,d\,e)$, so $k_1 = 1, m_1 = 2$ and $k_2 = 1, m_2 = 3$:
    \[
    \frac{5!}{(1! \cdot 2^1)(1! \cdot 3^1)} = \frac{5 \times 4 \times 3 \times 2 \times 1}{2 \times 3} = 20.
    \]
\end{enumerate}
\end{example}


\newpage
% =============================================================================
% PART II: INTRODUCTION TO RING THEORY
% =============================================================================
\part*{Part II: Introduction to Ring Theory}
\addcontentsline{toc}{part}{Part II: Introduction to Ring Theory}

\section{Introduction to Ring Theory: Definitions and Axioms}

\begin{definition}[Ring]
A algebraic system $(R, +, \cdot)$ consisting of a non-empty set $R$ equipped with two internal binary operations, denoted `$+$' (addition) and `$\cdot$' (multiplication), is called a \textbf{ring} if it satisfies the following three conditions:
\begin{enumerate}
    \item \textbf{$(R, +)$ is an abelian group:}
    \begin{itemize}
        \item \emph{Closure:} $a + b \in R$ for all $a, b \in R$.
        \item \emph{Associativity:} $(a + b) + c = a + (b + c)$ for all $a, b, c \in R$.
        \item \emph{Identity:} There exists an element $0 \in R$ such that $a + 0 = 0 + a = a$ for all $a \in R$.
        \item \emph{Inverse:} For each $a \in R$, there exists $-a \in R$ such that $a + (-a) = (-a) + a = 0$.
        \item \emph{Commutativity:} $a + b = b + a$ for all $a, b \in R$.
    \end{itemize}
    
    \item \textbf{$(R, \cdot)$ is a semi-group:}
    \begin{itemize}
        \item \emph{Closure:} $a \cdot b \in R$ for all $a, b \in R$.
        \item \emph{Associativity:} $(a \cdot b) \cdot c = a \cdot (b \cdot c)$ for all $a, b, c \in R$.
    \end{itemize}
    
    \item \textbf{Distributive Laws hold:} Multiplication distributes over addition from both left and right:
    \begin{align*}
    a \cdot (b + c) &= a \cdot b + a \cdot c \quad \forall a, b, c \in R, \\
    (a + b) \cdot c &= a \cdot c + b \cdot c \quad \forall a, b, c \in R.
    \end{align*}
\end{enumerate}
\end{definition}

\subsection{Terminology and Ring Properties}

\begin{enumerate}
    \item The operation `$+$' is usually referred to as \textbf{addition}, and `$\cdot$' is usually referred to as \textbf{multiplication}.
    
    \item The identity element of the abelian group $(R, +)$ is called the \textbf{zero element} of the ring and is denoted as $0$.
    
    \item The identity element with respect to multiplication `$\cdot$' (if it exists) is known as the \textbf{identity element of the ring} or \textbf{unity}, and is denoted as $1$.
    
    \item If the multiplication operation `$\cdot$' is commutative, i.e., $a \cdot b = b \cdot a$ for all $a, b \in R$, then $R$ is called a \textbf{commutative ring}.
    
    \item If the multiplicative identity exists in $R$, then $R$ is called a \textbf{ring with identity} or \textbf{ring with unity}.
\end{enumerate}

\section{Examples of Rings}

\subsection{Commutative Rings with Unity}

\begin{example}[Standard Number Systems]
The sets of numbers with standard addition and multiplication:
\[
(\mathbb{Z}, +, \cdot), \quad (\mathbb{Q}, +, \cdot), \quad (\mathbb{R}, +, \cdot), \quad (\mathbb{C}, +, \cdot)
\]
are all commutative rings with unity ($1$).
\end{example}

\begin{example}[Integers Modulo $n$]
For any positive integer $n \ge 2$, the set of congruence classes:
\[
(\mathbb{Z}_n, +, \cdot)
\]
under addition and multiplication modulo $n$ forms a commutative ring with unity ($[1]$).
\end{example}

\begin{example}[Quadratic Integer Ring]
Let $p$ be a prime. Define the set:
\[
\mathbb{Z}[\sqrt{p}] = \{a + b\sqrt{p} \mid a, b \in \mathbb{Z}\}.
\]
Then $\mathbb{Z}[\sqrt{p}]$ is a commutative ring with unity ($1 = 1 + 0\sqrt{p}$).
\end{example}

\begin{example}[Ring of Gaussian Integers]
The set of Gaussian integers:
\[
\mathbb{Z}[i] = \{a + bi \mid a, b \in \mathbb{Z}, \; i^2 = -1\}
\]
is a commutative ring with unity ($1 = 1 + 0i$).
\end{example}

\subsection{Non-Commutative Rings with Unity}

\begin{example}[Real Matrix Ring $M_n(\mathbb{R})$]
Let $M_n(\mathbb{R})$ denote the set of all $n \times n$ matrices with real entries:
\begin{enumerate}[label=(\roman*)]
    \item $(M_n(\mathbb{R}), +)$ is an abelian group with zero element the zero matrix $O_{n \times n}$.
    \item $(M_n(\mathbb{R}), \cdot)$ is a semi-group under matrix multiplication (associativity of matrix multiplication).
    \item The distributive laws hold:
    \begin{align*}
    A \cdot (B + C) &= AB + AC, \\
    (A + B) \cdot C &= AC + BC.
    \end{align*}
    \item The identity matrix $I_n$ satisfies $A I_n = I_n A = A$ for all $A \in M_n(\mathbb{R})$.
    \item Matrix multiplication is not commutative for $n \ge 2$ (in general, $AB \neq BA$).
\end{enumerate}
Therefore, $(M_n(\mathbb{R}), +, \cdot)$ is a \textbf{non-commutative ring with unity}.
\end{example}

\begin{example}[Integer Matrix Ring $M_n(\mathbb{Z})$]
Similarly, let $M_n(\mathbb{Z})$ denote the set of all $n \times n$ matrices with integer entries. Then $(M_n(\mathbb{Z}), +, \cdot)$ is also a \textbf{non-commutative ring with unity}.
\end{example}

\section{Special Matrix Ring: Identity and Inverse Computations}

\begin{example}[A Commutative Matrix Ring with Non-Standard Unity]
Consider the set:
\[
A = \left\{ \begin{bmatrix} x & x \\ x & x \end{bmatrix} \;\middle|\; x \in \mathbb{R} \right\}
\]
under standard matrix addition and multiplication.
\end{example}

\subsection{Verification of Ring Axioms}

\begin{enumerate}
    \item \textbf{$(A, +)$ is an abelian group:}
    For any $X = \begin{bmatrix} x & x \\ x & x \end{bmatrix}$ and $Y = \begin{bmatrix} y & y \\ y & y \end{bmatrix} \in A$:
    \[
    X + Y = \begin{bmatrix} x+y & x+y \\ x+y & x+y \end{bmatrix} \in A.
    \]
    Addition is associative and commutative since addition in $\mathbb{R}$ is. The zero element is:
    \[
    \mathbf{0} = \begin{bmatrix} 0 & 0 \\ 0 & 0 \end{bmatrix} \in A,
    \]
    and the additive inverse of $X$ is $-X = \begin{bmatrix} -x & -x \\ -x & -x \end{bmatrix} \in A$.
    
    \item \textbf{Commutativity of multiplication in $A$:}
    Let $X, Y \in A$ with $X = \begin{bmatrix} x & x \\ x & x \end{bmatrix}$ and $Y = \begin{bmatrix} y & y \\ y & y \end{bmatrix}$. Then:
    \begin{align*}
    \begin{bmatrix} x & x \\ x & x \end{bmatrix} \begin{bmatrix} y & y \\ y & y \end{bmatrix} &= \begin{bmatrix} x \cdot y + x \cdot y & x \cdot y + x \cdot y \\ x \cdot y + x \cdot y & x \cdot y + x \cdot y \end{bmatrix} \\
    &= \begin{bmatrix} 2xy & 2xy \\ 2xy & 2xy \end{bmatrix} \\
    &= \begin{bmatrix} 2yx & 2yx \\ 2yx & 2yx \end{bmatrix} \\
    &= \begin{bmatrix} y & y \\ y & y \end{bmatrix} \begin{bmatrix} x & x \\ x & x \end{bmatrix}.
    \end{align*}
    Therefore, multiplication is commutative in $A$.
    
    \item \textbf{Distributivity:}
    Matrix multiplication is always distributive over matrix addition.
\end{enumerate}
Thus, $(A, +, \cdot)$ is a \textbf{commutative ring}.

\subsection{Finding the Unity Element}

Let $E \in A$ be the multiplicative identity (unity) of $A$, so that:
\[
E = \begin{bmatrix} e & e \\ e & e \end{bmatrix}.
\]
Then $E$ must satisfy $X E = X$ for all $X \in A$:
\[
\begin{bmatrix} x & x \\ x & x \end{bmatrix} \begin{bmatrix} e & e \\ e & e \end{bmatrix} = \begin{bmatrix} x & x \\ x & x \end{bmatrix}.
\]
Performing the multiplication:
\[
\begin{bmatrix} 2xe & 2xe \\ 2xe & 2xe \end{bmatrix} = \begin{bmatrix} x & x \\ x & x \end{bmatrix}.
\]
Equating corresponding entries:
\[
2xe = x \implies 2e = 1 \implies e = \frac{1}{2}.
\]
Therefore, the identity element of the ring $A$ is:
\[
E = \begin{bmatrix} 1/2 & 1/2 \\ 1/2 & 1/2 \end{bmatrix}.
\]
\begin{remark}
Notice that the unity element $E \in A$ is \emph{not} the standard $2 \times 2$ identity matrix $I_2 = \begin{bmatrix} 1 & 0 \\ 0 & 1 \end{bmatrix}$ (which does not belong to $A$), yet it functions as the exact multiplicative identity on $A$.
\end{remark}

\subsection{Computation of Multiplicative Inverses}

For $(A \setminus \{\mathbf{0}\}, \cdot)$ to form an abelian group (making $A$ a field), each non-zero element must possess a multiplicative inverse. Let:
\[
Y = X^{-1} = \begin{bmatrix} y & y \\ y & y \end{bmatrix}.
\]
Then $X Y = E$:
\[
\begin{bmatrix} x & x \\ x & x \end{bmatrix} \begin{bmatrix} y & y \\ y & y \end{bmatrix} = \begin{bmatrix} 1/2 & 1/2 \\ 1/2 & 1/2 \end{bmatrix}.
\]
This gives:
\[
\begin{bmatrix} 2xy & 2xy \\ 2xy & 2xy \end{bmatrix} = \begin{bmatrix} 1/2 & 1/2 \\ 1/2 & 1/2 \end{bmatrix},
\]
which implies:
\[
2xy = \frac{1}{2} \implies y = \frac{1}{4x} \quad (x \neq 0).
\]
Therefore, for any non-zero element $X \in A$, its multiplicative inverse is:
\[
X^{-1} = \begin{bmatrix} \dfrac{1}{4x} & \dfrac{1}{4x} \\[6pt] \dfrac{1}{4x} & \dfrac{1}{4x} \end{bmatrix}.
\]
Consequently, $(A \setminus \{\mathbf{0}\}, \cdot)$ is an abelian group, and $A$ is isomorphic as a field to $\mathbb{R}$.


\end{document}
